
\documentclass{article}
\usepackage{colortbl}
\usepackage{makecell}
\usepackage{multirow}
\usepackage{supertabular}

\begin{document}

\newcounter{utterance}

\centering \large Interaction Transcript for game `imagegame', experiment `random\_grids', episode 9 with reka{-}flash{-}3.1{-}t0.0.
\vspace{24pt}

{ \footnotesize  \setcounter{utterance}{1}
\setlength{\tabcolsep}{0pt}
\begin{supertabular}{c@{$\;$}|p{.15\linewidth}@{}p{.15\linewidth}p{.15\linewidth}p{.15\linewidth}p{.15\linewidth}p{.15\linewidth}}
    \# & $\;$A & \multicolumn{4}{c}{Game Master} & $\;\:$B\\
    \hline

    \theutterance \stepcounter{utterance}  
    & & \multicolumn{4}{p{0.6\linewidth}}{
        \cellcolor[rgb]{0.9,0.9,0.9}{
            \makecell[{{p{\linewidth}}}]{
                \texttt{\tiny{[P1$\langle$GM]}}
                \texttt{Let us do the following task. The goal is to fill an empty grid that looks like this:} \\
\\ 
\texttt{▢ ▢ ▢ ▢ ▢} \\
\texttt{▢ ▢ ▢ ▢ ▢} \\
\texttt{▢ ▢ ▢ ▢ ▢} \\
\texttt{▢ ▢ ▢ ▢ ▢} \\
\texttt{▢ ▢ ▢ ▢ ▢} \\
\\ 
\texttt{A filled grid below is 5 by 5 and can look like this:} \\
\\ 
\texttt{▢ ▢ ▢ ▢ ▢} \\
\texttt{▢ ▢ E ▢ ▢} \\
\texttt{▢ ▢ ▢ ▢ ▢} \\
\texttt{▢ ▢ ▢ ▢ ▢} \\
\texttt{X X X X X} \\
\\ 
\texttt{I want you to describe this grid to me, step by step. You don't need to describe the empty squares, which are denoted with "▢". Only describe the location of letters in the grid. Then you wait for me to say "What is your next command?", and then you continue with the next step. Take the size of the grid into consideration while giving commands. You can give only one command at each turn but it can refer to single or multiple cells, rows, columns or even the whole grid at once. When you have described everything, you say "Command: DONE". Only give the command and omit everything else.} \\
\\ 
\texttt{For the filled grid above, here are the example steps.} \\
\\ 
\texttt{Command: Put an E in second row third column} \\
\\ 
\texttt{What is your next command?} \\
\texttt{Command: Fill the last row with X} \\
\\ 
\texttt{What is your next command?} \\
\texttt{Command: DONE} \\
\\ 
\texttt{Ok. Please do this for the following example, which is a 5 by 5 grid.} \\
\texttt{I ▢ ▢ I ▢} \\
\texttt{▢ ▢ ▢ ▢ I} \\
\texttt{I ▢ I I ▢} \\
\texttt{▢ ▢ ▢ ▢ I} \\
\texttt{▢ ▢ ▢ ▢ ▢} \\
\texttt{What is your next command?} \\
            }
        }
    }
    & \\ \\

    \theutterance \stepcounter{utterance}  
    & \multicolumn{4}{p{0.6\linewidth}}{
        \cellcolor[rgb]{0.9,0.9,0.9}{
            \makecell[{{p{\linewidth}}}]{
                \texttt{\tiny{[P1$\rangle$GM]}}
                \texttt{<reasoning>} \\
\texttt{Let's see. The user provided a 5x5 grid that's partially filled with some letters and empty squares. They want me to describe the steps to fill the grid by specifying commands, similar to the example they gave. The example had commands like putting an E in a specific cell and filling a row with Xs. So, I need to figure out the positions of the existing letters and then determine the commands needed to place them step by step.} \\
\\ 
\texttt{First, let me visualize the given grid. The user wrote:} \\
\\ 
\texttt{Row 1: I ▢ ▢ I ▢} \\
\texttt{Row 2: ▢ ▢ ▢ ▢ I} \\
\texttt{Row 3: I ▢ I I ▢} \\
\texttt{Row 4: ▢ ▢ ▢ ▢ I} \\
\texttt{Row 5: ▢ ▢ ▢ ▢ ▢} \\
\\ 
\texttt{Wait, but the user mentioned it's a 5x5 grid. Let me check each row to make sure. Each row should have 5 elements. Let's count:} \\
\\ 
\texttt{Row 1: I, ▢, ▢, I, ▢ → 5 elements. Correct.} \\
\texttt{Row 2: ▢, ▢, ▢, ▢, I → 5 elements. Correct.} \\
\texttt{Row 3: I, ▢, I, I, ▢ → 5 elements. Correct.} \\
\texttt{Row 4: ▢, ▢, ▢, ▢, I → 5 elements. Correct.} \\
\texttt{Row 5: ▢, ▢, ▢, ▢, ▢ → 5 elements. Correct.} \\
\\ 
\texttt{So the grid has some I's placed already. The task is to describe the steps to place these I's in the grid. The user wants me to give commands one by one, each referring to the placement of letters (probably I's here) in specific cells, rows, columns, or the whole grid. Then, after all commands, say DONE.} \\
\\ 
\texttt{Looking at the example, the first command was "Put an E in second row third column". Then "Fill the last row with X". So the commands can be specific cells or entire rows/columns. Since the example used both, maybe here I can do similar.} \\
\\ 
\texttt{Looking at the given grid, the I's are in the following positions:} \\
\\ 
\texttt{Row 1: columns 1 and 4 (assuming columns are 1{-}based)} \\
\texttt{Row 2: column 5} \\
\texttt{Row} \\
            }
        }
    }
    & & \\ \\

    \theutterance \stepcounter{utterance}  
    & & & \multicolumn{2}{p{0.3\linewidth}}{
        \cellcolor[rgb]{0.9,0.9,0.9}{
            \makecell[{{p{\linewidth}}}]{
                \texttt{\tiny{[GM$|$GM]}}
                \texttt{Invalid instruction format} \\
            }
        }
    }
    & & \\ \\

\end{supertabular}
}

\end{document}
